%==============================================================
%  lecons/analyse/riemann.tex — Intégrale de Riemann
%==============================================================

\lecontitle{Intégrale de Riemann}{Analyse réelle — Construction de l'intégrale}

%=============================================================
%  § 1 — DÉFINITION
%=============================================================
\secnum{navyblue}{1}{Subdivisions, sommes de Darboux et définition}

\begin{tcolorbox}[thmbox]
  \textbf{Cadre.}\enspace%
  \index{intégrale de Riemann}
  Soit $f:[a,b]\to\mathbb{R}$ \textbf{bornée} sur un segment.

  \medskip
  \textbf{Subdivision.}\enspace\index{subdivision}
  Une \emph{subdivision} de $[a,b]$ est une famille finie ordonnée
  $\sigma = (x_0, x_1, \ldots, x_n)$ avec
  \[
    a = x_0 < x_1 < \cdots < x_n = b.
  \]
  Son \emph{pas} est $|\sigma| = \max_{i}(x_i - x_{i-1})$.

  \medskip
  \textbf{Sommes de Darboux.}\enspace\index{sommes de Darboux}
  Pour une subdivision $\sigma$, on pose
  $m_i = \inf_{[x_{i-1},x_i]} f$ et $M_i = \sup_{[x_{i-1},x_i]} f$, puis
  \[
    s(f,\sigma) = \sum_{i=1}^{n} m_i\,(x_i - x_{i-1}),
    \qquad
    S(f,\sigma) = \sum_{i=1}^{n} M_i\,(x_i - x_{i-1}).
  \]
  Les \emph{intégrales de Darboux} sont
  \[
    \underline{I}(f) = \sup_\sigma s(f,\sigma),
    \qquad
    \overline{I}(f) = \inf_\sigma S(f,\sigma).
  \]
  On a toujours $\underline{I}(f) \leq \overline{I}(f)$.

  \medskip
  \textbf{Définition (intégrale de Riemann).}\enspace%
  \index{Riemann-intégrable}
  La fonction $f$ est dite \emph{Riemann-intégrable} sur $[a,b]$ si
  $\underline{I}(f) = \overline{I}(f)$. Cette valeur commune est notée
  \[
    \int_a^b f(t)\,\mathrm{d}t.
  \]

  \medskip
  \textbf{Théorème (critère de Riemann).}\enspace%
  $f$ est Riemann-intégrable sur $[a,b]$ si et seulement si
  \[
    \forall\varepsilon > 0,\ \exists\,\sigma\ \text{subdivision telle que}\
    S(f,\sigma) - s(f,\sigma) < \varepsilon.
  \]

  \medskip
  \textbf{Théorème (formulation par sommes de Riemann).}\enspace%
  \index{sommes de Riemann}
  $f$ est Riemann-intégrable si et seulement si il existe $I\in\mathbb{R}$ tel
  que, pour toute suite $(\sigma_n)$ de subdivisions de pas $|\sigma_n|\to 0$
  et tout choix $\xi_i^{(n)} \in [x_{i-1}^{(n)}, x_i^{(n)}]$,
  \[
    R(f,\sigma_n,\xi^{(n)}) \,=\, \sum_{i} f\bigl(\xi_i^{(n)}\bigr)\bigl(x_i^{(n)} - x_{i-1}^{(n)}\bigr)
    \xrightarrow[n\to+\infty]{} I = \int_a^b f.
  \]
\end{tcolorbox}

\begin{significationintuitive}
On approche l'aire sous la courbe par des rectangles. La somme inférieure
$s(f,\sigma)$ sous-estime l'aire (rectangles « par défaut »), la somme
supérieure $S(f,\sigma)$ la sur-estime (rectangles « par excès »).
Si les deux encadrements peuvent être rendus arbitrairement proches en
raffinant la subdivision, l'aire est bien définie : c'est l'intégrale.
L'idée géométrique remonte à Archimède (méthode d'exhaustion) ; la
formalisation rigoureuse est due à Riemann (1854).
\end{significationintuitive}

\decsep{}

%=============================================================
%  § 2 — PROPRIÉTÉS FONDAMENTALES
%=============================================================
\secnum{navyblue}{2}{Propriétés de l'intégrale}

\begin{tcolorbox}[thmbox]
  Soient $f, g : [a,b] \to \mathbb{R}$ Riemann-intégrables et
  $\lambda, \mu \in \mathbb{R}$.
  \begin{enumerate}
    \item \textbf{Linéarité :}
          $\lambda f + \mu g$ est intégrable et
          $\displaystyle\int_a^b (\lambda f + \mu g) = \lambda\!\int_a^b f + \mu\!\int_a^b g$.

    \item \textbf{Positivité :}
          $f \geq 0 \Rightarrow \int_a^b f \geq 0$.
          Si de plus $f$ est continue et $\int_a^b f = 0$, alors $f \equiv 0$.

    \item \textbf{Croissance :}
          $f \leq g \Rightarrow \int_a^b f \leq \int_a^b g$.

    \item \textbf{Inégalité triangulaire :}
          $|f|$ est intégrable et
          $\displaystyle\left|\int_a^b f\right| \leq \int_a^b |f| \leq (b-a)\,\|f\|_{\infty,[a,b]}$.

    \item \textbf{Relation de Chasles :}\index{Chasles (relation de)}
          pour $c\in[a,b]$, $f$ intégrable sur $[a,c]$ et $[c,b]$ et
          $\displaystyle\int_a^b f = \int_a^c f + \int_c^b f$.

    \item \textbf{Produit :} $f\cdot g$ est intégrable.

    \item \textbf{Formule de la moyenne.}\enspace\index{formule de la moyenne}
          Si $f$ est continue, il existe $c\in[a,b]$ tel que
          $\displaystyle\int_a^b f = (b-a)\,f(c)$.
  \end{enumerate}
\end{tcolorbox}

\rubriq{Classes de fonctions Riemann-intégrables}

\begin{tcolorbox}[thmbox]
  Sur tout segment $[a,b]$, sont Riemann-intégrables :
  \begin{itemize}
    \item les fonctions \textbf{continues}~;
    \item les fonctions \textbf{monotones} (même si elles présentent des
          discontinuités, du moment qu'elles sont bornées)~;
    \item les fonctions \textbf{continues par morceaux}~;
    \item plus généralement (\textbf{critère de Lebesgue}\index{critère de Lebesgue}) :
          $f$ bornée est Riemann-intégrable si et seulement si son ensemble
          de points de discontinuité est de \emph{mesure nulle}.
  \end{itemize}
\end{tcolorbox}

\decsep{}

%=============================================================
%  § 3 — DÉMONSTRATIONS
%=============================================================
\secnum{navyblue}{3}{Démonstrations}

\rubriq{Critère de Riemann}

\begin{mdframed}[style=proofbar]
  \textit{$(\Rightarrow)$}\enspace%
  Si $f$ est intégrable, $\underline{I}=\overline{I}=I$. Par définition du
  sup et de l'inf, pour $\varepsilon>0$, il existe deux subdivisions
  $\sigma_1, \sigma_2$ telles que $s(f,\sigma_1) > I - \varepsilon/2$ et
  $S(f,\sigma_2) < I + \varepsilon/2$. La subdivision plus fine $\sigma$
  réunion de $\sigma_1$ et $\sigma_2$ vérifie
  $s(f,\sigma)\geq s(f,\sigma_1)$ et $S(f,\sigma)\leq S(f,\sigma_2)$
  (le raffinement augmente $s$ et diminue $S$), donc
  $S(f,\sigma) - s(f,\sigma) < \varepsilon$.

  \medskip
  \textit{$(\Leftarrow)$}\enspace%
  Inversement, $s(f,\sigma) \leq \underline{I} \leq \overline{I} \leq S(f,\sigma)$
  pour tout $\sigma$. Donc $\overline{I} - \underline{I} \leq S(f,\sigma) - s(f,\sigma) < \varepsilon$
  pour tout $\varepsilon>0$ ; on a $\underline{I} = \overline{I}$.
  \hfill$\blacksquare$
\end{mdframed}

\rubriq{Continuité $\Rightarrow$ intégrabilité}

\begin{mdframed}[style=proofbar]
  Soit $f$ continue sur $[a,b]$. Par le théorème de Heine, $f$ est
  \emph{uniformément continue} : pour $\varepsilon>0$, il existe $\delta>0$
  tel que $|x-y| < \delta \Rightarrow |f(x)-f(y)| < \varepsilon/(b-a)$.

  \medskip
  Soit $\sigma$ une subdivision de pas $|\sigma| < \delta$. Sur chaque
  intervalle $[x_{i-1},x_i]$, $M_i - m_i \leq \varepsilon/(b-a)$ (atteint sur
  un compact par continuité), donc
  \[
    S(f,\sigma) - s(f,\sigma) = \sum_i (M_i - m_i)(x_i - x_{i-1})
    \leq \tfrac{\varepsilon}{b-a}\sum_i (x_i - x_{i-1}) = \varepsilon.
  \]
  Le critère de Riemann conclut.\hfill$\blacksquare$
\end{mdframed}

\rubriq{Monotonie $\Rightarrow$ intégrabilité}

\begin{mdframed}[style=proofbar]
  Soit $f$ croissante sur $[a,b]$ (donc bornée, de bornes $f(a)$ et $f(b)$).
  Sur une subdivision $\sigma$ de pas $|\sigma|$, $m_i = f(x_{i-1})$ et
  $M_i = f(x_i)$ par croissance. En majorant chaque facteur
  $(x_i - x_{i-1}) \leq |\sigma|$ et en utilisant le télescopage
  $\sum_i \bigl(f(x_i) - f(x_{i-1})\bigr) = f(b) - f(a)$ (somme télescopique :
  les bornes intermédiaires $f(x_i)$ s'annulent deux à deux), on obtient
  \[
    S(f,\sigma) - s(f,\sigma) = \sum_i \bigl(f(x_i) - f(x_{i-1})\bigr)(x_i - x_{i-1})
    \leq |\sigma|\sum_i \bigl(f(x_i) - f(x_{i-1})\bigr) = |\sigma|\cdot\bigl(f(b) - f(a)\bigr).
  \]
  Pour $|\sigma|$ assez petit, c'est $< \varepsilon$. Le critère
  s'applique.\hfill$\blacksquare$
\end{mdframed}

\rubriq{Équivalence Darboux $\Leftrightarrow$ sommes de Riemann}

\begin{mdframed}[style=proofbar]
  \textit{Encadrement.}\enspace%
  Pour toute subdivision pointée $(\sigma, \xi)$ avec
  $\xi_i\in[x_{i-1},x_i]$,
  \[
    s(f,\sigma) \leq R(f,\sigma,\xi) = \sum_i f(\xi_i)(x_i - x_{i-1}) \leq S(f,\sigma)
  \]
  car $m_i \leq f(\xi_i) \leq M_i$.

  \medskip
  \textit{$(\Rightarrow)$}\enspace%
  Si $f$ est intégrable, le critère de Riemann donne
  $S(f,\sigma_n)-s(f,\sigma_n)\to 0$ pour $|\sigma_n|\to 0$. L'encadrement
  ci-dessus force $R(f,\sigma_n,\xi^{(n)})\to I$.

  \medskip
  \textit{$(\Leftarrow)$}\enspace%
  En choisissant les $\xi_i$ pour approcher $m_i$ ou $M_i$ à $\varepsilon$
  près (sur compact, supremum et infimum sont atteints ou approchés), on
  voit que $s(f,\sigma)$ et $S(f,\sigma)$ peuvent être obtenus comme
  limites de sommes de Riemann. Si toutes ces sommes convergent vers
  $I$, alors $\underline{I} = \overline{I} = I$.\hfill$\blacksquare$
\end{mdframed}

\rubriq{Formule de la moyenne}

\begin{mdframed}[style=proofbar]
  Pour $f$ continue, posons $m = \min_{[a,b]} f$ et $M = \max_{[a,b]} f$.
  Par croissance, $m(b-a) \leq \int_a^b f \leq M(b-a)$, donc
  $\mu := \dfrac{1}{b-a}\int_a^b f \in [m, M]$. Le théorème des valeurs
  intermédiaires\index{théorème des valeurs intermédiaires} donne
  $c\in[a,b]$ avec $f(c) = \mu$.\hfill$\blacksquare$
\end{mdframed}

\decsep{}

%=============================================================
%  § 4 — EXEMPLES, CONTRE-EXEMPLES, EXERCICES
%=============================================================
\secnum{exgreen}{4}{Exemples, contre-exemples et exercices}

\rubriq{Exemples}

\begin{mdframed}[style=exbar]
  \textbf{1. Fonction constante.}\enspace%
  Si $f \equiv c$, alors $s(f,\sigma) = S(f,\sigma) = c(b-a)$ pour toute
  subdivision : $\int_a^b c\,\mathrm{d}t = c(b-a)$.

  \smallskip\noindent
  \textbf{2. Identité par sommes de Riemann.}\enspace%
  $\displaystyle\int_0^1 t\,\mathrm{d}t = \lim_{n\to+\infty}\frac{1}{n}\sum_{k=1}^{n}\frac{k}{n}
   = \lim_{n\to+\infty}\frac{n+1}{2n} = \frac{1}{2}$.

  \smallskip\noindent
  \textbf{3. Carré par sommes de Riemann.}\enspace%
  Avec la subdivision régulière $x_k = k/n$,
  $\displaystyle\int_0^1 t^2\,\mathrm{d}t = \lim_{n\to+\infty}\frac{1}{n}\sum_{k=1}^{n}\frac{k^2}{n^2}
   = \lim_{n\to+\infty}\frac{(n+1)(2n+1)}{6n^2} = \frac{1}{3}$.

  \smallskip\noindent
  \textbf{4. Reconnaissance d'une somme de Riemann.}\enspace%
  $\displaystyle\lim_{n\to+\infty}\sum_{k=1}^{n}\frac{n}{n^2 + k^2}
   = \int_0^1 \frac{\mathrm{d}t}{1+t^2} = \frac{\pi}{4}$.

  \smallskip\noindent
  \textbf{5. Fonction monotone discontinue intégrable.}\enspace%
  Posons $f(t) = \sum_{n\geq 1} 2^{-n}\,\mathbf{1}_{[1/n,\,1]}(t)$ et $f(0)=0$.
  Cette fonction est croissante, \emph{bornée} par $\sum_{n\geq 1}2^{-n} = 1$,
  et présente un saut de hauteur $2^{-n}$ en chaque point $1/n$ : elle possède
  donc une infinité dénombrable de discontinuités tout en restant
  Riemann-intégrable (toute fonction monotone bornée l'est). On prendra garde
  qu'une fonction monotone \emph{non bornée}, comme $t\mapsto\lfloor 1/t\rfloor$
  au voisinage de $0$, n'est en revanche pas Riemann-intégrable.
\end{mdframed}

\vspace{6pt}
\rubriq{Contre-exemples et pièges}

\begin{mdframed}[style=cexbar]
  \textbf{Fonction de Dirichlet.}\enspace\index{fonction de Dirichlet}
  $\mathbf{1}_{\mathbb{Q}}:[0,1]\to\{0,1\}$ : sur tout sous-intervalle non
  trivial, $m_i = 0$ et $M_i = 1$, donc $s(f,\sigma) = 0$ et $S(f,\sigma) = 1$
  pour toute $\sigma$. Donc $\underline{I} = 0 \neq 1 = \overline{I}$ :
  cette fonction n'est \emph{pas} Riemann-intégrable. Elle est Lebesgue-intégrable
  d'intégrale $0$ (l'ensemble $\mathbb{Q}\cap[0,1]$ est de mesure nulle).

  \smallskip\noindent
  \textbf{Fonction non bornée.}\enspace%
  $f(t) = 1/\sqrt{t}$ sur $]0,1]$ n'est pas bornée, donc \emph{pas}
  Riemann-intégrable sur $[0,1]$ au sens strict. On peut cependant définir
  une \emph{intégrale impropre} $\int_0^1 t^{-1/2}\,\mathrm{d}t = 2$
  par passage à la limite $\int_\varepsilon^1$.

  \smallskip\noindent
  \textbf{La continuité ponctuelle ne suffit pas pour l'intégrabilité.}\enspace%
  La fonction de Thomae $f(p/q) = 1/q$ (fraction irréductible) et $f(x) = 0$
  pour $x$ irrationnel est continue en chaque irrationnel mais discontinue
  sur $\mathbb{Q}$. Elle reste Riemann-intégrable d'intégrale $0$
  (ensemble de discontinuités dénombrable, donc de mesure nulle).
\end{mdframed}

\vspace{6pt}
\exolabel

\begin{mdframed}[style=exobar]
  \begin{enumerate}
    \item Calculer $\displaystyle\int_0^1 t^2\,\mathrm{d}t$ par sommes
          de Riemann ($\sum_{k=1}^n k^2 = n(n+1)(2n+1)/6$).

    \item Montrer que $\displaystyle\lim_{n\to+\infty}\frac{1}{n}\sum_{k=1}^{n}\sqrt{\frac{k}{n}} = \int_0^1\sqrt{t}\,\mathrm{d}t = \frac{2}{3}$.

    \item Calculer
          $\displaystyle\lim_{n\to+\infty}\prod_{k=1}^{n}\!\left(1+\frac{k}{n}\right)^{1/n}$
          en reconnaissant $\exp$ d'une somme de Riemann.

    \item Démontrer que toute fonction continue par morceaux sur $[a,b]$
          est Riemann-intégrable. (\textit{Indication}~: linéarité +
          continuité sur chaque morceau, puis Chasles.)

    \item Soit $f:[0,1]\to\mathbb{R}$ continue. Montrer que
          $\displaystyle\lim_{n\to+\infty}\int_0^1 f(t^n)\,\mathrm{d}t = f(0)$.

    \item Construire une fonction Riemann-intégrable sur $[0,1]$ dont
          l'ensemble des discontinuités est exactement $\{1/n \mid n\geq 1\}\cup\{0\}$.

    \item (\textit{Critère de Lebesgue, sens facile}) Si $f:[a,b]\to\mathbb{R}$
          bornée admet un ensemble fini de discontinuités, montrer que $f$
          est Riemann-intégrable par découpage et étude au voisinage de chaque
          point de discontinuité.

    \item Montrer que la moyenne $\frac{1}{b-a}\int_a^b f$ d'une fonction
          continue est limite des moyennes discrètes
          $\frac{1}{n}\sum_{k=1}^n f(a + k(b-a)/n)$.
  \end{enumerate}
\end{mdframed}

\leconfooter{B.\ Riemann (1854) — G.\ Darboux (1875) — H.\ Lebesgue (1904)}
